\section*{Author notation and preliminaries}
We use $[m]$ to represent $\{1,2,3,\ldots,m\}$ and $\binom m{\le r}$ to denote $\binom m0+\cdots+\binom mr$. We denote by $\binom{[m]}r$ and $\binom{[m]}{\le r}$ the sets $\{S\subseteq[m]:|S|=r\}$ and $\{S\subseteq[m]:|S|\le r\}$ respectively.
We shall abuse notation and use $\mathrm{RM}(m,r)$ to denote the vector space of all $m$-variate multivariate polynomials in $\mathbb F_2[x_1,\ldots,x_m]$ of degree at most $r$. Throughout the paper, the parameter $m$ will be unchanged and we will avoid mentioning it for brevity. Let $n=2^m$.
For parameters $m,r$, define the matrix $E(m,r)$ as the $\binom m{\le r}\times2^m$-matrix whose columns are indexed by elements in $\mathbb F_2^m$ and rows by multlinear $m$-variate monomials of degree at most $r$ and whose $z$-th column is the vector consisting of the evaluation of all multilinear $m$-variate degree $\le r$ monomials on $z$.
For a polynomial $P\in\mathrm{RM}(m,r)$, we use $\mathrm{wt}(P)$ to denote the fractional Hamming weight:
\[\mathrm{wt}(P):=\frac{|\{z\in\mathbb F_2^m:P(z)\ne0\}|}{2^m}.\]
We shall say that a polynomial $P\in\mathrm{RM}(m,r)$ is $\eta$-biased if $|\mathrm{wt}(P)-1/2|\ge\eta$. We shall also use $\mathrm{RM}_\eta(m,r)$ to denote the set of all $\eta$-biased polynomials in $\mathrm{RM}(m,r)$. For a real number $\alpha\in(0,1)$, we denote by $\mathrm{WtDist}_{m,r}(\alpha)$ the number of polynomials of (fractional) weight at most $\alpha$ in $\mathrm{RM}(m,r)$, i.e., $\mathrm{WtDist}_{m,r}(\alpha):=|\{P\in\mathrm{RM}(m,r):\mathrm{wt}(p)\le\alpha\}|$. All complexity notations used in the paper are with respect to $m$ as the growing parameter.
\paragraph{Theorem 2.1 (Sberlo and Shpilka [13]: bounds for low-weight codewords).} For any $\ell\ge1$, we have
\[\mathrm{WtDist}_{m,r}(2^{-\ell})\le\operatorname{exp}_2\left(O(m^4)+17(c_\gamma\ell+d_\gamma)\gamma^{\ell-1}\binom m{\le r}\right)\]
where $c_\gamma=1/(1-\gamma)$ and $d_\gamma=(2-\gamma)/(1-\gamma)^2$. In particular, if $\gamma\le1/2$, we have
\[\mathrm{WtDist}_{m,r}(2^{-\ell})\le\operatorname{exp}_2\left(O(m^4)+O(\ell\gamma^{\ell-1})\binom m{\le r}\right).\tag{1}\]
\paragraph{Theorem 2.2 (Sberlo and Shpilka [13]: bounds for medium-weight codewords).} Assume that $\gamma\in(0,(1/2)-\Omega(1))$. For a positive integer $\ell$ such that $\ell/m$ upper bounded by a small enough constant,\footnote{[13, Theorem 1.3] only guarantees this for $\ell=o(m)$ but [13, Remark 1.1] after the theorem statement says that in this setting it holds for $\ell=\Omega(m)$.} we have
\[\mathrm{WtDist}_{m,r}\left(\frac12-2^{-\ell}\right)\le\operatorname{exp}_2\left(O(m^4)+(1-2^{-c(\gamma,\ell)})\binom m{\le r}\right),\tag{2}\]
where $c(\gamma,\ell)=O(\max\{\gamma^2\ell,\gamma\})$.
\section*{Degree-$r$ vanishing spaces and closures}
We first define the notion of a vanishing space. This is similar to the notion of vanishing ideals in basic algebraic geometry but we refer to them as vanishing space instead to stress that we are studying them as a vector space and not an ideal.
\paragraph{Definition 3.1 (degree-$r$ vanishing spaces).} For a set $Z\subseteq\mathbb F_2^m$, we use $I_r(Z)$ to denote the degree-$r$ vanishing space defined as
\[I_r(Z):=\{P\in\mathrm{RM}(m,r):P(z)=0\text{ for all }z\in Z\}.\]
Related to the vanishing ideals is also the notion of the degree-$r$ closure (similar to the Zariski closure in standard algebraic geometry but restricted to the setting of $\mathbb F_2^m$), which is the set of points on which every polynomial in the degree-$r$ vanishing space vanishes.
\paragraph{Definition 3.2 (degree-$r$ closure).} For a set $Z\subseteq\mathbb F_2^m$, we use $\mathrm{Closure}_r(Z)$ to denote the degree-$r$ closure of $Z$ defined as
\[\mathrm{Closure}_r(Z):=\{u\in\mathbb F_2^m:P(u)=0\text{ for all }P\in I_r(Z)\}.\]
The above notion can be equivalently defined as the set of all $u\in\mathbb F_2^m$ such that column of $E(m,r)$ indexed by $u$ is in the span of columns of $E(m,r)$ indexed by $z\in Z$. That is,
\[\mathrm{Closure}_r(Z)=\{u\in\mathbb F_2^m:u^{(r)}\in\operatorname{span}\{z^{(r)}:z\in Z\}\},\]
where $z^{(r)}$ denotes the $\binom m{\le r}$-dimension vector of evaluations of all $m$-variate degree at most $r$ monomials at the point $z$. For any set $Z\subseteq\mathbb F_2^m$ of size $k$, note that $I_r(Z)$ is a vector space of dimension at least $\binom m{\le r}-k$, as each constraint $P(z)=0$ adds one homogeneous linear constraint on the ambient space $\mathrm{RM}(m,r)$. In fact, it is easy to see that $I_r(Z)$ has rank $\binom m{\le r}-k$ if and only if each point of $Z$ is not in the closure of the previous points.
\paragraph{Observation 3.3 (vanishing ideals of minimal rank).} Let $Z=\{z_1,\ldots,z_k\}\subseteq\mathbb F_2^m$. Then,
\[\dim I_r(Z)=\binom m{\le r}-k\iff(\forall i=1,\ldots,k-1:z_i\notin\mathrm{Closure}_r(\{z_1,\ldots,z_{i-1}\})).\]
